
\item
\emph{Les parties @@M0@@ de @@M1@@ correspondent biunivoquement aux familles @@M2@@ de classes de faisceaux inversibles sur les fibres de @@M3@@.}
\item \emph{Une famille @@M4@@ de @@M5@@ est quasi-compacte si et seulement si la famille @@M6@@ correspondante est limit\'ee.}
\end{enumerate}

En effet, (i) r\'esulte de (3.1)\,; (ii), du fait suivant~: si @@M7@@ est un point de @@M8@@ et @@M9@@ est une extension alg\'ebriquement close de @@M10@@, alors les @@M11@@-points de @@M12@@ correspondent biunivoquement aux classes d'isomorphisme de faisceaux inversibles sur @@M13@@.

Quant \`a (iii), supposons que la partie @@M14@@ de @@M15@@ soit quasi-compacte. Alors @@M16@@ est contenue dans un ouvert @@M17@@ de type fini, et l'inclusion @@M18@@ correspond \`a un faisceau inversible @@M19@@ sur un @@M20@@, o\`u @@M21@@ est un morphisme surjectif, (plat) de type fini. Enfin @@M22@@ est de type fini, donc @@M23@@ limite @@M24@@.

R\'eciproquement supposons que @@M25@@ soit limit\'ee par un faisceau coh\'erent @@M26@@ sur un @@M27@@, avec @@M28@@ de type fini sur @@M29@@. On peut supposer @@M30@@ inversible (1.13 (i)) afin que @@M31@@ d\'efinisse un morphisme @@M32@@. L'image de @@M33@@ contient @@M34@@ et elle est quasi-compacte, donc noeth\'erienne par (i). Par suite @@M35@@ est quasi-compact et on a (iii).

Gr\^ace \`a 3.2, la d\'efinition habituelle des morphismes de type fini entre deux sch\'emas de Picard se g\'en\'eralise en la suivante~:

\textbf{D\'efinition 3.3.} \emph{Soient @@M36@@ et @@M37@@ propres sur @@M38@@ noeth\'erien. Un morphisme entre les foncteurs de Picard}
@@M39@@
